%
% Rationalisierung
%
\subsection{Rationalisierung
\label{buch:intalg:wurzel:subsection:rationalisierung}}
Die Darstellung~\eqref{buch:algint:wurzel:eqn:funktion} der Funktion kann
weiter vereinfacht werden, mit potenzieller Vereinfachung der Berechnung
der Stammfunktion.
Dazu wird \eqref{buch:algint:wurzel:eqn:funktion} mit
$q_0-q_1\vartheta$ zu
\begin{align}
f
&=
\frac{p_0+p_1\vartheta}{q_0+q_1\vartheta}
\cdot
\frac{q_0-q_1\vartheta}{q_0-q_1\vartheta}
\notag
\\
&=
\frac{
p_0q_0 +(p_0q_1+p_1q_0)\vartheta + p_1q_1\vartheta^2
}{
q_0^2-q_1^2\vartheta^2
}
\notag
\\
&=
\frac{
p_0q_0+p_1q_1(ax^2+bx+c)
}{
q_0^2-q_1^2(ax^2+bx+c)
}
+
\frac{
p_0q_1+p_1q_0
}{
q_0^2-q_1^2(ax^2+bx+c)
}
\vartheta
\notag
\\
&=
p + q\vartheta
\label{buch:intalg:wurzel:eqn:linkombtheta}
\end{align}
erweitert.
Daraus liest man ab, dass der Funktionenkörper sogar als
\[
\mathbb{Q}(x,\vartheta)
=
\{
p+q\vartheta
\mid
p,q\in\mathbb{Q}(x)
\}
\]
geschrieben werden kann.

Die Relation $\vartheta^2-(ax^2+bx+c)=0$ für das algebraische Element
$\vartheta$ bedeutet auch, dass
\[
\vartheta
=
\frac{ax^2+bx+c}{\vartheta}
\]
und damit, dass der Funktionenkörper auch aus den Linearkombinationen
\begin{equation}
\mathbb{Q}(x,\vartheta)
=
\biggl\{
p+q\frac{1}{\vartheta}
\bigg|
p,q\in\mathbb{Q}(x)
\biggr\}
\label{buch:intalg:wurzel:eqn:linkombthetainverse}
\end{equation}
von $1$ und $1/\vartheta$ mit Koeffizienten in $\mathbb{Q}(x)$ besteht.
